\section{The residual decomposition and its parametric reading}\label{sec:parametric}

Repeatedly extracting a maximum-ratio set and continuing on the residual is not
a device of this project. It is the decomposition of \citet{Sidney1975} for
precedence-constrained single-machine scheduling, developed in the sequencing
literature by \citet{Lawler1978} and surveyed by \citet{CorreaSchulz2005};
\citet{MargotEtAl2003} give the version closest to the ratio reading used here:
an in-depth study of an LP relaxation of a precedence-constrained problem whose
algorithmic core is a parametric extension of Sidney's decomposition, obtained
by essentially solving a parametric minimum-cut problem.

The vocabulary of this document follows \citet{DoseEtAl2026}, the companion work
of this project, which studies the same canonical object on a precedence graph
that is a directed forest. What is called a \emph{canonical block} here is a
\emph{closure layer} there; the canonical decomposition is their \emph{optimal
sequence of closure layers}; the ratio $\lambda_r$ of a block is a
\emph{breakpoint} of $u(\lambda)$; and the per-vertex bound this project
computes is their \emph{threshold} of a vertex --- the same quantity, not an
analogy. Their scalar $\lambda$, parametrising the weights as $p_i - \lambda
w_i$, is the $\rho$ of the ratio oracle here, which is also the reduced cost
$b_i = p_i - \rho w_i$ of \cref{sec:lg}. The two works do not overlap on
content: they give aggregation algorithms specialised to forests, reaching
$\mathcal{O}(n\log n)$, whereas the precedence graph here is an arbitrary DAG
and the method is the simplex. The sequence produced is the same object, so the
two are directly comparable on instances that happen to be forests --- and on
those instances their bound is the one to beat.
The same object appears from the flow side as the sequence of nested minimum
cuts of a parametric network: \citet{EisnerSeverance1976} give the
divide-and-conquer scheme that computes all breakpoints with a number of
maximum-flow computations proportional to their number, and
\citet{GalloEtAl1989} give the parametric maximum-flow algorithm that computes
the whole path in the time of a single flow computation for monotone
parametrisations. \citet{Topkis1978} supplies the lattice/submodularity view
that makes the nesting structural rather than incidental.
\citet{PicardSmith2004} apply exactly this to open-pit contours, and
\citet{Eppstein2018} gives the modern complexity treatment of the parametric
closure problem.

This project implements the Eisner--Severance scheme as its \texttt{parametric}
backend, and its experimental report is explicit that the implementation is a
divide-and-conquer decomposition with independent maximum flows, without the
amortised flow reuse of \citet{GalloEtAl1989}.

\paragraph{What this costs the claim.} The existence of the canonical
decomposition, its nesting, the strictly decreasing ratios, the interpolation
between two consecutive prefixes that yields the LP value at a given capacity,
and the solution of the whole path by parametric flow, are all established. The
project's mathematical contract already presents them as such. A claim of
novelty at this level would be simply wrong, and the earlier drafts were right
to drop it.

\citet{NemeschEtAl2025} survey exact and approximation algorithms for
linear-parametric optimisation and place the closure case inside the general
picture; it is the reference to consult before asserting that any parametric
technique used here is unusual.

\paragraph{What survives here.} Only one thing, and it is small: the convention
that resolves ties. When several closures attain the maximum ratio, the
canonical block is their union, which on the flow side is the \emph{maximal}
source side of a minimum cut rather than the minimal one
\citep{PicardQueyranne1980,PicardQueyranne1982}. This project implements both
conventions and verifies that they agree; that is careful engineering on a known
structure, not a new result.
